\section{充分统计量保留了哪些参数信息}
\label{sec:11-Mathematical-Statistics/03-Sufficient-Statistics-and-Information/01}

一条线上服务每天写下四亿行请求日志。存储按季度扩容，账单已经压过了这些日志本身带来的收益，于是有人提议改成只落盘几个聚合量：请求数、成功数、耗时之和、耗时平方和。省钱这件事没人反对，反对的理由是另一句话——原始日志删掉就再也回不来，以后要重算某个结论怎么办。

这句反对成不成立，取决于以后要问什么。如果问的是本周成功率有没有下滑，那么在一个说得清楚的模型下，请求数与成功数已经把该带的都带上了，把全量日志重新扫一遍也翻不出新东西。如果问的是失败是不是集中在某几分钟，聚合量就答不上来，因为它把时间顺序抹平了。同一次压缩，对一个问题免费，对另一个问题致命。

先看最小的版本。抛硬币 5 次，两台机器分别记下

\[
(1,0,1,1,0)\quad\text{和}\quad(0,1,1,0,1).
\]

两串数据不同。但在 \(X_i\overset{\iid}{\sim}\Bernoulli(\theta)\) 模型下，似然只依赖成功总数 \(S=\sum_i X_i\)：

\[
p_\theta(x)=\theta^{S}(1-\theta)^{n-S},
\]

顺序压根没进公式。两台机器对 \(\theta\) 的推断一字不差。

\begin{center}
\begin{tikzpicture}[font=\small,>=stealth]
  \node[draw,rounded corners,minimum width=3.6cm,minimum height=0.9cm,align=center] (X) at (0,3.5) {原始样本 \(X\)：全量日志};
  \node[draw,rounded corners,minimum width=3.4cm,minimum height=0.9cm,align=center] (T1) at (-3.3,1.85) {\(T_1=\) 随机留 \(10\%\)};
  \node[draw,rounded corners,minimum width=3.4cm,minimum height=0.9cm,align=center] (T2) at (3.3,1.85) {\(T_2=\bigl(n,\ \sum_i X_i\bigr)\)};
  \draw[->,black!60] (X) -- (T1);
  \draw[->,black!60] (X) -- (T2);
  \fill[blue!25] (-4.9,0.55) rectangle (-3.2,1.0);
  \draw[black!45] (-4.9,0.55) rectangle (-1.7,1.0);
  \fill[blue!55!black] (1.7,0.55) rectangle (4.9,1.0);
  \draw[black!45] (1.7,0.55) rectangle (4.9,1.0);
  \node[font=\footnotesize,anchor=north,align=center] at (-3.3,0.35) {关于 \(\theta\) 的信息\\少了一截};
  \node[font=\footnotesize,anchor=north,align=center] at (3.3,0.35) {关于 \(\theta\) 的信息\\一点没少};
  \node[font=\footnotesize,black!65] at (0,-1.25) {灰框长度 \(=\) 全量日志能提供的信息};
\end{tikzpicture}

\emph{两条压缩路径，左边掉了一截信息，右边恰好没掉；充分性说的就是右边这种情形。}
\end{center}

图里两个方框都是压缩，都把 \(n\) 条记录换成短得多的摘要。差别在下面那两根条：左边那根短了一截，右边那根顶到灰框末端。本节要做的是把``顶到末端''写成一个可以验证的条件，并且说明为什么这种压缩存在得如此频繁——它并非侥幸，而是模型形状带来的结果。

\subsection{给定 \(T\) 之后，原始数据对参数不再说话}

统计量 \(T(X)\) 对参数 \(\theta\) 充分，指的是给定 \(T(X)=t\) 后，完整样本 \(X\) 的条件分布不再依赖 \(\theta\)。

这句定义的操作含义值得慢读一遍。把 \(T\) 的取值告诉你以后，剩余的随机性——观测排在哪个位置、哪一条恰好偏高、同一天内部怎样打散——仍然存在，但它们的分布是一张与 \(\theta\) 无关的表。既然那张表里没有 \(\theta\)，看再多遍也调整不了你对 \(\theta\) 的判断。日志那件事因此有了答案：只要 \(T\) 对你关心的参数充分，删掉原始记录不会削弱针对该参数的推断能力。

麻烦在于这个定义不好直接验证。要检查``条件分布不依赖参数''，得先把条件分布写出来，而条件分布往往比原始密度还难写。

\subsection{因子分解定理把充分性变成一次代数观察}

Neyman--Fisher 因子分解定理绕开了条件分布。在共同支配测度这类常规条件下，

\[
T(X)\ \text{充分}\quad\Longleftrightarrow\quad p_\theta(x)=g_\theta\bigl(T(x)\bigr)\,h(x),
\]

其中 \(g_\theta\) 只通过 \(T(x)\) 接触数据，\(h(x)\) 与 \(\theta\) 无关。

验证变成了一次代数观察：写下联合密度，看参数是从哪个入口进来的。Bernoulli 样本里 \(p_\theta(x)=\theta^{\sum_i x_i}(1-\theta)^{n-\sum_i x_i}\)，取 \(T(x)=\sum_i x_i\)、\(g_\theta(t)=\theta^t(1-\theta)^{n-t}\)、\(h(x)=1\)，分解立刻成立。成功总数对 \(\theta\) 充分，图里右边那根条到此有了依据。

条件对一遍。因子分解要求密度对同一个测度成立，且 \(h\) 完全不含 \(\theta\)，包括藏在示性函数里的那部分——这一点在本节末尾会栽一次跟头。至于 \(g_\theta\) 的形状则不受限制，它可以任意难看，只要参数不越过 \(T\) 这道门。

模型一变，结论就跟着变。如果硬币的概率随时间漂移，或者相邻两次抛掷相关，顺序不再是无关细节，它开始携带关于漂移速度或相关系数的信息，\(\sum_i x_i\) 也就不再充分。充分性始终是统计量、模型、参数三者之间的关系，不是统计量单独拥有的属性。线上日志的场景里，这意味着聚合方案与分析口径必须一起冻结：改了要问的问题，可能就得改存的东西。

\subsection{一般的压缩会掉信息，充分统计量是恰好不掉的那种}

\hyperref[sec:07-Information-Theory/02-Joint-Information/04]{``数据处理不等式与充分表示''}给过一条一般判断：沿 Markov 链 \(\theta\to X\to T\)，有 \(I(\theta;T)\le I(\theta;X)\)。任何只作用在 \(X\) 上的确定性变换或加噪都不会创造关于 \(\theta\) 的信息，最多原样保留。这条不等式适用于图里的两条路径，它解释了左边那根条为什么短——随机抽 \(10\%\) 属于典型的有损处理。

有意思的是等号成立的条件。\(I(\theta;T)=I(\theta;X)\) 当且仅当 \(\theta\to T\to X\) 也构成 Markov 链，也就是给定 \(T\) 后 \(X\) 与 \(\theta\) 条件独立——这正是上面那条定义。充分性于是有了一个信息论读法：不等式对所有压缩都成立，充分统计量是把它取成等号的那些压缩。频率学派框架里 \(\theta\) 不是随机变量，互信息写不出来，但只要临时给 \(\theta\) 一个先验，两套语言就对得上；而因子分解形式无需借道先验，这也是它在数理统计里更常用的原因。

两处术语容易串。信息论里的``充分表示''常指对标签 \(Y\) 保留预测所需的信息，这里的充分性针对的是模型参数 \(\theta\)。二者结构相同，压缩的目标对象不同。

\subsection{哪些参数未知，决定了什么摘要算充分}

取 \(X_i\overset{\iid}{\sim}N(\mu,\sigma^2)\)，两个参数都未知。对数似然

\[
\ell(\mu,\sigma^2)=-\frac n2\log(2\pi\sigma^2)-\frac{1}{2\sigma^2}\Bigl(\sum_i X_i^2-2\mu\sum_i X_i+n\mu^2\Bigr)
\]

里，数据只通过两个数进场。因子分解给出

\[
T(X)=\Bigl(\sum_i X_i,\ \sum_i X_i^2\Bigr),
\]

也可以等价地写成 \((\bar X,S^2)\)，二者互为可逆变换，携带的信息相同。请求日志里存耗时之和与耗时平方和，存的就是它。

把 \(\sigma^2\) 改成已知，摘要立刻缩短。此时含 \(\sum_i X_i^2\) 的那一项与 \(\mu\) 无关，可以整块划进 \(h(x)\)，\(\bar X\) 一个数就对 \(\mu\) 充分。哪些量被声明为未知，直接改变了什么样的摘要够用。这句话落到存储方案上是有后果的：今天认定方差已知而只存均值，明天想估方差就没有原料了。

\subsection{摘要长度不随样本量增长，才是聚合方案可行的原因}

上面两个例子有个共同点：无论 \(n\) 是五还是四亿，充分统计量都是固定的两三个数。这不是巧合。密度能写成

\[
p_\theta(x)=h(x)\exp\bigl\{\eta(\theta)^{\trans}T(x)-A(\theta)\bigr\}
\]

的模型称为指数族，Bernoulli、Poisson、正态、Gamma 都在其中；因子分解对它自动成立，且 \(T\) 的维数由参数个数定，与样本量无关。反过来的结论同样成立：在支持集不随参数变化的常规条件下，只有指数族才有维数固定的充分统计量。指数族的其他好处见\hyperref[sec:13-Machine-Learning/06-Probabilistic-Models/06]{``指数族把归一化、矩与似然连接起来''}。

工程上的后果是直接的。维数固定意味着摘要可以增量更新：新到一批数据，把它的计数与和累加上去即可，无须回看历史，也无须为``再攒一个月''预留额外字段。流式统计、分片汇总、按天分桶再合并，靠的都是这条性质。跳出指数族则往往没有这种运气——非参数模型里可能根本不存在固定长度的充分统计量，随着样本增加，你要保留的摘要也一直变长，聚合方案就退化成了压缩率不高的存储方案。

\subsection{充分不等于低维，也不等于估计得好}

一个顺手的误读是把充分当成``压缩成低维摘要''。原始数据本身 \(T(X)=X\) 永远充分，它一个字节都没省。充分性只承诺不掉信息，压缩多少是另一回事，两件事得分开验收。

另一头，充分统计量不是估计量。\(T\) 可能是向量、矩阵甚至函数，你仍然要决定怎样由它构造点估计、区间或检验，而充分性对这些规则的偏差与方差不作任何承诺。它给的保证只有一句：手里有了 \(T\)，翻回原始数据不会带来关于 \(\theta\) 的新信息。

不过充分性确实能兑换成一次估计改进。设 \(U(X)\) 是任意一个估计量，\(T\) 充分，令 \(\tilde U(T)=\E_\theta[U(X)\given T]\)。因为 \(T\) 充分，条件分布不含 \(\theta\)，这个条件期望是你算得出来的规则。由重期望公式 \(\E[\tilde U]=\E[U]\)，两者偏差相同；再由方差分解

\[
\Var(U)=\Var\bigl(\E[U\given T]\bigr)+\E\bigl[\Var(U\given T)\bigr]
\]

右端两项非负，第一项即 \(\Var(\tilde U)\)，于是 \(\Var(\tilde U)\le\Var(U)\)，进而 \(\MSE_\theta(\tilde U)\le\MSE_\theta(U)\)。等号成立当且仅当给定 \(T\) 后 \(U\) 几乎处处是常数，也就是 \(U\) 原本就没沾 \(T\) 之外的随机细节。这是 Rao--Blackwell 定理，它的用法很具体：估计量若还依赖排列顺序或别的与参数无关的波动，对充分统计量取一次条件期望就能把那部分抹平，方差只降不升。

\subsection{支持集随参数移动时，因子分解要小心}

Uniform\((0,\theta)\) 是标准反例。似然

\[
p_\theta(x)=\theta^{-n}\,\ind_{\max_i x_i\le\theta}
\]

看上去可以把 \(\theta^{-n}\) 放进 \(g_\theta\)、示性函数放进 \(h(x)\)，但示性函数里写着 \(\theta\)，它没有资格待在 \(h\) 里。正确的做法是令 \(g_\theta(t)=\theta^{-n}\ind_{t\le\theta}\)、\(h(x)=1\)，取 \(T=\max_i X_i\)。样本最大值仍然充分，只是通往结论的路径与前面几例不同：这里携带信息的是支持集的边界，而不是密度的形状。上一小节关于指数族的那句``反过来也成立''，排除的正是这类模型。

还有一处技术性提醒。连续模型里``给定 \(T=t\)''常常是零概率事件，本科课程中密度相除的做法在此不够用，需要正规条件分布，或者干脆用因子分解形式绕开条件化。定义层面的严格性藏在测度论里，不影响使用，但影响你对``条件分布''这四个字能推多远的判断。

现在换一个角度追问。充分性回答的是哪些数据细节多余，它是一道是非题：信息要么全在 \(T\) 里，要么不全在。可是日志留下来之后，真正决定推断精度的是另一件事——这批数据对 \(\theta\) 到底说得有多死。同样是一百万条记录，估计一个接近 \(0.5\) 的成功率和估计一个接近 \(0.999\) 的成功率，难度差得很远，而两者的充分统计量是同一个。要把``说得有多死''量化成一个数，需要看似然函数在真值附近的形状，这是下一节的事。
